\documentclass[11pt,a4paper,final,oneside]{article}

% -----------------------------------------------------------------------------
% Encoding and typography
% -----------------------------------------------------------------------------
\usepackage[T1]{fontenc}
\usepackage[utf8]{inputenc}
\usepackage[english]{babel}
\usepackage{lmodern}
\usepackage{microtype}
\usepackage{times}

% -----------------------------------------------------------------------------
% Page layout and document utilities
% -----------------------------------------------------------------------------
\usepackage[a4paper,top=2.0cm,bottom=2.0cm,left=2.5cm,right=2.5cm]{geometry}
\usepackage{graphicx}
\usepackage{xcolor}
\usepackage{array}
\usepackage{tabularx}
\usepackage{fancyhdr}
\usepackage{lastpage}
\usepackage[hidelinks]{hyperref}
\usepackage{enumitem}
\usepackage{seqsplit}

\definecolor{IST_blue}{rgb}{0,0.1922,0.8863}

% -----------------------------------------------------------------------------
% Header and footer
% -----------------------------------------------------------------------------
\pagestyle{fancy}
\fancyhf{}
\renewcommand{\headrulewidth}{0pt}
\fancyfoot[C]{\footnotesize CSF -- Tutorial I -- Page~\thepage~of~\pageref{LastPage}}
\setlength{\headheight}{14pt}

% -----------------------------------------------------------------------------
% Paragraph formatting
% -----------------------------------------------------------------------------
\setlength{\parindent}{0pt}
\setlength{\parskip}{0.48em}

\begin{document}
\thispagestyle{empty}

% -----------------------------------------------------------------------------
% Title page
% -----------------------------------------------------------------------------
\begin{center}

\begin{minipage}[c]{0.18\textwidth}
  \centering
  \includegraphics[width=\linewidth]{figures/NewLogoIST.jpg}
\end{minipage}
\hfill
\begin{minipage}[c]{0.76\textwidth}
  \centering
  {\LARGE\bfseries\sffamily INSTITUTO SUPERIOR TÉCNICO\par}
  \vspace{0.20cm}
  {\large\sffamily Departamento de Engenharia Informática\par}
\end{minipage}

\vspace{2.3cm}

{\huge\bfseries\sffamily Forensics Cyber-Security\par}
\vspace{0.35cm}
{\large MEIC, METI\par}

\vfill

{\LARGE\bfseries Digital Forensics Report\par}
\vspace{0.40cm}
{\Large\bfseries Tutorial 1 -- Introduction to File Forensics\par}

\vspace{1.2cm}

% TODO: Replace with your group details.
\renewcommand{\arraystretch}{1.35}
\begin{tabularx}{0.82\textwidth}{>{\bfseries}l X r}
Group: & \multicolumn{2}{l}{Group number} \\
\hline
Student 1: & Name 1 & Number 1 \\
Student 2: & Name 2 & Number 2 \\
Student 3: & Name 3 & Number 3 \\
\end{tabularx}

\vfill

{\large 2026/2027\par}

\end{center}

\clearpage

% -----------------------------------------------------------------------------
% Acquired artifacts
% -----------------------------------------------------------------------------
\section{Acquired artifacts}
\label{sec:acquired_artifacts}

\begin{center}
\renewcommand{\arraystretch}{1.4}
\begin{tabularx}{\textwidth}{
>{\raggedright\arraybackslash}p{0.25\textwidth}
>{\centering\arraybackslash}p{0.14\textwidth}
>{\raggedright\arraybackslash}X}
\hline
\textbf{Name} & \textbf{Type} & \textbf{Hash / String} \\
\hline
\texttt{sd-transfer-test.jpg} & SHA-256 &
{\scriptsize\texttt{\seqsplit{de3c0a760aab64cf5f1ef8aa86736e5e09fd93c743771e7899b89ac4dc048bc8}}} \\
SSH key file password & String & \texttt{SSH key pwd: pfg2EPH2eyb@kyp7bjv} \\
\texttt{schedule.pdf} & SHA-256 &
{\scriptsize\texttt{\seqsplit{97b7d4120cc69bb33b64e6666a62428d00e7bbd9fe31a8f854bac999947a0d13}}} \\
\texttt{secret.pdf} & SHA-256 &
{\scriptsize\texttt{\seqsplit{3ab4c4b4d7163e4e01a16a71df5096ce0996d3becabf43a133d6901b03297a76}}} \\
\hline
\end{tabularx}
\end{center}

% -----------------------------------------------------------------------------
% Report of all findings
% -----------------------------------------------------------------------------
\section{Report of all findings}
\label{sec:report-of-all-findings}

We have discovered 4 secrets, three of which are files and one that is a password.

The first secret we found was a photo of a Rubik's cube with an SD card on top, which we found by analyzing the file called \texttt{sd-transfer-test}. The first thing we noticed was that the file did not have an extension, meaning that it was not obvious what type of file it was. Therefore, we started by using the following command to obtain the first few bytes of the file in search of its magic numbers:

\begin{quote}
\texttt{hexdump sd-transfer-test | head -n 5}
\end{quote}

The output of this command showed us that the first 4 bytes of the file were ``ff d8 ff db''. Having obtained this sequence of bytes, we then proceeded to check if they corresponded to any file type's magic numbers, by searching for it in the ``List of file signatures'' Wikipedia page. There we discovered that these bytes correspond to magic numbers of the JPEG file type. Finally, we simply added the ``.jpeg'' extension to the file and attempted to open it.

The second secret we found was a password for an SSH key file, which we found within the contents of the \texttt{cat.jpg} file. We started our analysis by opening the image to verify if there were any visible indicators of tampering, such as through the LSB method. As we did not find anything suspicious in the image itself, we decided to continue our analysis by searching for any relevant strings that might have been hidden within the contents of the file. To do this we executed the following command:

\begin{quote}
\texttt{strings -n 8 cat.jpg}
\end{quote}

It is important to note that we decided to use the \texttt{-n 8} argument to reduce the amount of randomly occurring strings outputted by the command, thus making it easier to search for relevant strings. The resulting output ended up being small enough that we immediately noticed the secret string hidden within the file: ``SSH key pwd: pfg2EPH2eyb@kyp7bjv''.

The third secret we found was a PDF file describing a recipe for a ``Secret French Dressing'', which we found by cracking the password of the \texttt{protected.zip} file. As expected, we started the analysis of the \texttt{protected.zip} file by attempting to open it and extract its contents, however, in attempting to do so we immediately discovered that it was protected with a password. This meant that to successfully extract the file that we could see was inside, named \texttt{secret-recipe.pdf}, we would have to somehow discover the password. Thankfully, password protected zip files are not hard to crack using bruteforce attacks. With that objective in mind we proceeded to install John the Ripper and to use it together with the rockyou wordlist to launch a bruteforce attack against the zip file. This was done through two commands, the first to extract the hash to crack from the zip file, and the other to truly launch the attack:

\begin{quote}
\texttt{zip2john protected.zip > hash.txt}\\
\texttt{john --wordlist=/usr/share/wordlists/rockyou.txt hash.txt}
\end{quote}

The bruteforce attack ended up being quite quick in finding the password, which was ``zaq1XSW@''. Having found the password of the zip file, we were finally able to extract the secret file, allowing us to get more context on the crime being committed.

The fourth and final secret that we found was a PDF file containing a schedule for security guards, which we found by analyzing the \texttt{broken.pdf} file. The first thing we attempted to do was open the file to verify what its contents were, however, that resulted in an error saying that the file is either not a PDF or corrupted. This meant that we should first verify what type of file it could be instead, which we did by running the following command:

\begin{quote}
\texttt{file broken.pdf}
\end{quote}

The output of the \texttt{file} command suggested that the file may not be a PDF but instead a text file. As this was suspicious we decided to further investigate the first bytes of the file, running the following command:

\begin{quote}
\texttt{xxd broken.pdf | head -n 1}
\end{quote}

By analyzing the first few ASCII characters of the file, we immediately theorized that the pattern and amount of \% characters together with the presence of the string PDF could suggest that the original PDF file had suffered some form of encoding. In an attempt to verify our theory, we decided to search online for an encoding that seemed to match the patterns we could observe. We quickly found out about URL encoding, which not only seemed to match what we saw, but made us realize that the ``\%25'', ``\%0D'' and ``\%0A'' found at the beginning of the file were the encoded forms of the characters \%, CR (carriage return) and LF (line feed), which together with the ``PDF'' in between, corresponded to the magic numbers of the PDF format.

With this in mind, we searched online for a tool that would allow us to decode the file. We ended up deciding to use the hURL command-line tool, which allowed us to run a simple command to obtain the secret file hidden within the \texttt{broken.pdf}:

\begin{quote}
\texttt{hURL -u -s --nocolor -f broken.pdf > schedule.pdf}
\end{quote}

% -----------------------------------------------------------------------------
% Analysis of relevant findings
% -----------------------------------------------------------------------------
\section{Analysis of relevant findings}
\label{sec:analysis-of-relevant-findings}

\subsection{Likely identity of the storage device's owner}

\textit{What can you conclude about the likely identity of the storage device's owner? Justify your answer with relevant findings.}

Following our examination of the documents, we proceeded to analyze the metadata of each file to uncover clues about the owner of the storage device. Fortunately, within the metadata, we discovered a file attributed to ``Eduardo Toca de Neve'' in the ``Author'' field. It seems probable that this file was deliberately transferred to the SD card as a means of testing the transfer method, considering its name. Furthermore, given that the file itself is a photo of an SD card, it's plausible to infer that it was captured by its owner.

\begin{figure}[htbp]
    \centering
    \includegraphics[width=0.43\textwidth]{figures/rubiks-sd-card.jpg}
    \caption{Photo of an SD card on top of a Rubik's cube.}
    \label{fig:rubiks-sd}
\end{figure}

However, the presence of this file amidst other discreet and sensitive files hints at possible carelessness or an intentional diversion by the suspect. Nevertheless, despite this ambiguity, Eduardo Toca de Neve emerges as the prime suspect based on the available evidence.

\subsection{Indications of illicit activity and interpretation of the data}

\textit{Based on the secrets you recovered, is there any indication of illicit activities involving the acquired SD card? If there's no direct evidence of any such activity, how would you interpret the data? Formulate a hypothesis regarding their purpose and justify it using the content of the recovered secrets.}

The files recovered from the SD card contain notably sensitive information, particularly the document housing the secret recipe. Not only is this file marked with a prominent ``CONFIDENTIAL'' stamp, but it also includes a note emphasizing the importance of not sharing the document.

\begin{figure}[htbp]
    \centering
    \includegraphics[width=0.92\textwidth]{figures/recipe-note.png}
    \caption{Note written on the secret recipe document.}
    \label{fig:recipe-note}
\end{figure}

In conjunction with this, the presence of a schedule detailing security guard shifts, with a specific shift highlighted, implies a connection to illicit activities. It suggests that the SD card may have been utilized in an illegal scheme to steal confidential documents from a highly secure location. Furthermore, the inclusion of a photo depicting the SD card concealed within a Rubik's cube strongly indicates a covert method of bypassing security checks.

\begin{figure}[htbp]
    \centering
    \includegraphics[width=0.93\textwidth]{figures/security-guard-shifts.png}
    \caption{Schedule of security guards with one shift highlighted.}
    \label{fig:guard-shifts}
\end{figure}

Therefore, it is reasonable to hypothesize that the purpose of this SD card was to facilitate the clandestine extraction of sensitive information, potentially for illicit gains. This interpretation is substantiated by the contents of the recovered secrets, which collectively point towards a covert operation involving theft and subterfuge.

\subsection{Recommendations for the subsequent course of action}

\textit{Given your discoveries, what would be your recommendations for the subsequent course of action?}

Based on our findings, we recommend the following steps to continue the investigation:

\begin{itemize}[leftmargin=1.6em]
    \item Investigate individuals that might be associated with the SD card, particularly Eduardo Toca de Neve, for possible involvement in illicit activities.
    \item Request access to security footage and seek witness accounts to corroborate suspicions of unauthorized access or theft.
    \item Request access to pertinent network captures and system logs within the National Salad Assembly's internal network.
    \item Upon obtaining further confirmation of Eduardo's or another person's involvement, it is advisable to request a warrant to search their residence for additional evidence.
    \item Coordinate with the authorities to initiate legal proceedings if criminal activity is confirmed.
\end{itemize}

% -----------------------------------------------------------------------------
% Use of Artificial Intelligence
% -----------------------------------------------------------------------------
\section{Use of Artificial Intelligence}
\label{sec:ai}

\textit{Did you use Artificial Intelligence (AI) tools in the preparation
of this report? If so, identify the tool(s) used and briefly state the
purpose(s) for which they were used.}

% TODO:
% Answer ``No'', or identify the AI tool(s) used and the purpose(s) of use.


% -----------------------------------------------------------------------------
% Appendices
% -----------------------------------------------------------------------------
\section{Appendices}
\label{sec:appendices}

% TODO:
% Attach relevant supporting evidence, screenshots, commands,
% tool outputs or other material useful to substantiate or
% reproduce the findings.

Feel free to attach appendices, e.g., displaying relevant evidence, timelines, etc.

\end{document}
